\section{Construction}\label{sec:tsps-construction}

The signing key of an underlying SPS scheme $\SPS$ is secret-shared among the
$n$ signers. Rather than reconstructing it, the signers use the MPC protocol of
\Cref{sec:tsps-mpc} to jointly compute shares of every value that the signing
algorithm derives from the key and from its own randomness.

\paragraph{Setup.} The signers agree in advance on $(\ell, t, n)$ and the
pairing groups, and abort if any call to $\algo{MPC}$ returns $\bot$.

\paragraph{Key generation and verification.} Key generation runs $\SPSgen$
under the MPC protocol. Verification is $\SPSverify$ unchanged, since it uses no
secret shares.

\begin{outline}
  \item both generic: general-purpose MPC gives a distributed key generation with
no trusted dealer; it runs once, so its cost hardly matters
  \item nothing for the MPC protocol to do in verification
\end{outline}

\begin{outline}
  \item evaluating $\SPSsign(\sk, \vec{\msg})$ under MPC with the message as one
more input puts every interactive call after the message is known
  \item instead split signing so all interactive calls happen before the message
is known and the online phase is local; this depends on the signing algorithm,
which we define first
\end{outline}

\paragraph{Decomposable signing.} A scheme $\SPS$ has decomposable signing if every
signature $\SPSsig \sample \SPSsign(\sk, \vec{\msg})$ on $\vec{\msg} = (\msg_1,
\ldots, \msg_\ell)$ consists of group elements $c_0$, which do not depend on the
message, and $z_1, \ldots, z_k$ of the following form:
%
\begin{equation*}
%
  z_j = c_j \prod_{i=1}^{\ell} \msg_i^{e_{j,i}},
%
\end{equation*}
%
where the group elements $c_0, c_1, \ldots, c_k$ and field elements $e_{j,i}$
are computed from $\sk$ and the signing randomness alone, using the interfaces
of $\algo{MPC}$ and linear operations.\rebekah{Defining this through the MPC
interfaces makes it vacuous, since any function can be computed with them. It
should be a property of $\SPS$ itself. The definition also allows $c_j$ and
$e_{j,i}$ to be linear in $\sk$, as in $\KPW$ where $e_{j,i}$ are key entries.}

\begin{outline}
  \item why this makes the online phase local: from sharings of $c_j$ and
$e_{j,i}$, each signer raises the public $\msg_i$ to its shares of $e_{j,i}$ and
multiplies by its share of $c_j$, a linear operation
  \item not every scheme has this form: e.g.\ message inside an inverse,
$\gen^{1/(x+\msg)}$, needs an interactive call once the message is known
\end{outline}

The construction $\TSPS$, applied to an SPS scheme $\SPS$, is given in
\Cref{prot:tsps-generic}.

\begin{protocolbox}{Threshold structure-preserving signatures}{tsps-generic}

  \begin{description}

    \item[$(\shr{\sk}, \pk) \sample \TSPSgen(\secparam, \ell, t, n)$ :] The
signers do the following:

      \begin{enumerate}

	\item Run $\SPSgen(\secparam, \ell)$ under $\algo{MPC}$: every field
element it samples is a call to $\algo{RandShare}$, and every product of secret
field elements it derives is a call to $\algo{Mul}$.

	\item Compute the shares of the group elements of $\pk$ and $\sk$
locally, by raising the generators to the shares of their discrete logarithms.

	\item Call $\algo{Open}$ on $\pk$ only, which broadcasts each signer's
share $\pk_i$ of it.

      \end{enumerate}

    \item[$\shr{\state} \sample \TSPSoffsign(\shr{\sk})$ :] The signers do the
following:

      \begin{enumerate}

	\item Sample the signing randomness with $\algo{RandShare}$.

	\item Compute sharings of $c_0, \ldots, c_k$ and of every $e_{j,i}$ with
$\algo{Mul}$, $\algo{Inv}$, $\algo{ExpGrp}$ and linear operations.

	\item Store $\shr{\state} = (\shr{c_0}, \ldots, \shr{c_k},
(\shr{e_{j,i}})_{j,i})$.

      \end{enumerate}

    \item[$\SPSsig_i \leftarrow \TSPSonsign(\sk_i, \state_i, \vec{\msg})$ :]
Each signer does the following, with no call to $\algo{MPC}$:

      \begin{enumerate}

	\item Compute the following for each $j \in [1,k]$:
	  %
	  \begin{equation*}
	  %
	    \shr{z_j} = \shr{c_j} \prod_{i=1}^{\ell} \msg_i^{\shr{e_{j,i}}}.
	  %
	  \end{equation*}

	\item Return its shares of $(c_0, z_1, \ldots, z_k)$.

      \end{enumerate}

    \item[$\SPSsig \leftarrow \TSPScombine(\{\SPSsig_i\}_{i \in Q})$ :]
Reconstruct each group element of $\SPSsig$ from the $t+1$ shares, and return
$\SPSsig$.

    \item[$\{0,1\} \leftarrow \TSPSverify(\pk, \vec{\msg}, \SPSsig)$ :] Return
$\SPSverify(\pk, \vec{\msg}, \SPSsig)$.

  \end{description}

\end{protocolbox}

\paragraph{Offline signing.} Each state is used for one message only. For both schemes we instantiate, two
signatures under the same state give a
relation between the two messages and the secret key, so the signers keep their
states in a queue and delete each one once it is used. The offline phase can be
run repeatedly during quiet periods to build up states for busy ones.

\paragraph{Online signing.} Online signing takes the message vector and
produces a signature share. By decomposable signing, the only operations needed
at this point are multiplication of shares and raising the public group elements
of $\vec{\msg}$ to shared exponents, both of which each signer performs locally.
No signer communicates with any other during this phase, and its cost is
comparable to that of an ordinary invocation of $\SPSsign$.

\paragraph{Combining.} Combining reconstructs the signature from the shares. It
is performed by whichever party is to hold the signature, and needs no
participation from the signers beyond their having published their shares. In
our benchmarks, every signer opens the signature, so each sends its share to all
others.\rebekah{Combining is defined as non-interactive by any holder, but bad
shares are only caught at the authenticated opening, and the benchmarks measure
every party opening. Say how a holder combining alone detects a bad share.}
\Cref{fig:tsps-seq} shows the steps and messages of signing, from the
offline phase to combining.

\begin{figure}[htbp] \centering \small
\begin{tikzpicture}[seq, x=2.7cm]
  \foreach \x/\n in {0/Signer $1$, 1/\ldots, 2/Signer $n$, 3/Holder} {
    \node[seqhead, text width=2.3cm] at (\x, 0) {\n};
  }
  \foreach \x in {0, 2, 3} \draw[seqlife] (\x, 0.6) -- (\x, 7.5);
  \draw[seqlife, dotted] (1, 0.6) -- (1, 7.5);
  \draw[{Stealth[length=1.6mm]}-{Stealth[length=1.6mm]}, thin] (0, 1.3) --
    node[seqlabel] {$\TSPSoffsign$ under $\algo{MPC}$} (2, 1.3);
  \node[seqloc] at (0, 2.1) {keep $\state_1$};
  \node[seqloc] at (2, 2.1) {keep $\state_n$};
  \draw[gray!60, dashed] (-0.4, 3.3) -- node[seqlabel, text=black] {$\vec{\msg}$ known} (3.4, 3.3);
  \node[seqloc] at (0, 4.0) {$\SPSsig_1 \leftarrow \TSPSonsign$};
  \node[seqloc] at (2, 4.0) {$\SPSsig_n \leftarrow \TSPSonsign$};
  \draw[seqmsg] (2, 5.5) -- node[seqlabel] {$\SPSsig_n$} (3, 5.5);
  \draw[seqmsg] (0, 6.4) -- node[seqlabel, pos=0.25] {$\SPSsig_1$} (3, 6.4);
  \node[seqloc] at (3, 7.2) {$\SPSsig \leftarrow \TSPScombine$};
\end{tikzpicture}
\caption[Threshold SPS schemes: signing]{Steps and messages of signing. The
signers run the offline phase together before the message is known, and each
then signs locally. The holder combines the signature shares of any $t+1$
signers.} \label{fig:tsps-seq} \end{figure}

\paragraph{Robustness and abort.} With an honest majority, the honest signers' shares alone are enough to combine,
so signing is robust: it succeeds whatever the corrupted signers do. With
additive sharing, every signer has to contribute, and corrupted signers can make
signing abort. If the MPC protocol has identifiable abort, meaning that an abort
also names a corrupted signer, such signers can be excluded.\rebekah{Robust
overstates it: the offline MPC can abort even with an honest majority, and
$\TSPScombine$ has to detect bad shares for the honest shares to be enough.}

\begin{outline}
  \item what we prove: thresholdising loses nothing; if $\SPS$ can't be forged
and the MPC reveals nothing beyond what it opens, $\TSPS$ can't be forged
  \item only need EUF-CMA of $\SPS$, not strong unforgeability: TS-UF-1 counts a
forgery by its message, so the reduction never has to tell two signatures on the
same message apart; asking for more would rule out rerandomisable schemes like
$\AGHO$
  \item from the MPC: privacy and security with abort against up to $t$
adaptively corrupted signers, the linear $(t+1, n)$ sharing of
\Cref{sec:tsps-mpc}
  \item from $\SPS$: decomposable signing, which is what lets the online phase
run without the MPC
\end{outline}

\begin{lemma}\label{lem:tsps-construction-correctness} Let $\SPS$ be a correct SPS
scheme with decomposable signing, and let the MPC protocol be secure with abort
over a linear $(t+1, n)$ threshold sharing. Then $\TSPS$ applied to $\SPS$ is
correct.
\end{lemma}

We prove \Cref{lem:tsps-construction-correctness} in \Cref{sec:tsps-proof-correctness}.


\begin{theorem}\label{thm:tsps-construction-security} Let $\SPS$ be an EUF-CMA
secure SPS scheme with decomposable signing, and let the MPC protocol be private
and secure with abort against up to $t$ adaptively corrupted signers, over a
linear $(t+1, n)$ threshold sharing. For every PPT adversary $\adv$ against the
adp-TS-UF-1 security of $\TSPS$ applied to $\SPS$, there is a PPT adversary
$\advB$ such that the following holds:
  %
  \begin{equation*}
  %
    \Adv_{\TSPS, \adv}^{\mathsf{adp\text{-}TS\text{-}UF\text{-}1}}(\secparam) \leq \binom{n}{t}
\cdot \Adv_{\SPS, \advB}^{\mathsf{euf\text{-}cma}}(\secparam).
  %
  \end{equation*}
  %
\end{theorem}

We prove \Cref{thm:tsps-construction-security} in \Cref{sec:tsps-proof-construction}.


\begin{outline}
  \item if $n$ is polynomial in $\secparam$ and $t$ or $n - t$ is constant, then
$\binom{n}{t} \leq n^{\min(t, n-t)}$ is polynomial, so $\TSPS$ is adp-TS-UF-1
secure
  \item why the factor disappears below: the reduction answers honest signers
lazily instead of guessing
\end{outline}

The factor $\binom{n}{t}$ disappears if the adversary fixes its corrupted signers
in advance.

\begin{lemma}\label{lem:tsps-static} Under the conditions of
\Cref{thm:tsps-construction-security}, with the MPC protocol secure against static
corruption only, for every PPT adversary $\adv$ that fixes
its corrupted signers before key generation and makes no corruption queries,
there is a PPT adversary $\advB$ such that the following holds for any $n$ and
$t$:
  %
  \begin{equation*}
  %
    \Adv_{\TSPS, \adv}^{\mathsf{adp\text{-}TS\text{-}UF\text{-}1}}(\secparam) \leq
\Adv_{\SPS, \advB}^{\mathsf{euf\text{-}cma}}(\secparam).
  %
  \end{equation*}
  %
\end{lemma}

We prove \Cref{lem:tsps-static} in \Cref{sec:tsps-static}.


\begin{outline}
  \item with an adaptively secure protocol: dishonest majority $t = n-1$ loses
only $\binom{n}{n-1} = n$; honest majority $t \approx n/2$ loses a factor
exponential in $n$, only covered for constant $n$
  \item where the factor comes from: the reduction guesses, before any
corruption, which signers the adversary will corrupt or ask for shares on its
forgery message
  \item removing it adaptively would need an equivocable sharing (simulator can
later explain it as a sharing of any value) or secure erasures (honest signers
reliably delete states); neither benchmarked protocol provides these
\end{outline}
